\chapter{Fazit}
\label{chap:fazit}

Diese Ausarbeitung hat drei Forschungsfragen zur Dynamik von Speedrunning-Weltrekorden auf \textit{speedrun.com} untersucht. Die zentralen Erkenntnisse lassen sich wie folgt zusammenfassen:

Die \textbf{Zeitreduktion} variiert zwischen Spielen erheblich (\textbf{9,4~\%–98,9~\%}) und wird stärker von spielspezifischen Faktoren wie Spiellänge und verfügbaren Glitches bestimmt als vom Genre (kein signifikanter Genreunterschied, $\mathbf{p = 0{,}483}$).

Der \textbf{Sättigungspunkt} ist in allen 17~Spielen statistisch nachweisbar: Der Chow-Test zeigt universelle Strukturbrüche, die typischerweise in der frühen Rekordphase auftreten. Die Verbesserungskurve lässt sich mit hoher Güte modellieren ($\mathbf{R^2 > 0{,}73}$), wobei kein einzelnes Modell für alle Spiele optimal ist.

Der \textbf{externe Einfluss} auf die Rekordentwicklung wurde in drei Teilfragen aufgeteilt. \textbf{F3c~(Lebensdauer)} ist vollständig beantwortet: Die Lebensdauer unterscheidet sich signifikant zwischen Genres ($\mathbf{H = 43{,}05}$; $\mathbf{p < 0{,}001}$). \textit{RPG}-Rekorde halten mit einem Median von \textbf{52~Tagen} deutlich länger als Rekorde anderer Genres (\textbf{6–15~Tage}); der Rückgang ab den 2010er~Jahren korreliert mit dem Wachstum der Speedrunning-Community.

\textbf{F3a~(Post-Breakthrough-Dynamik)} quantifiziert erstmals, ob und wie stark ein identifizierbares Durchbruchereignis die Verbesserungsrate einer Community verändert. Das Flattening Ratio $\phi$ operationalisiert diesen Effekt reproduzierbar auf Basis öffentlich zugänglicher WR-Daten.

\placeholder{Kernergebnis von F3a eintragen: Wie viele Spiele zeigen ein Plateau ($\phi < 0{,}5$) vs.\ Beschleunigung ($\phi > 1{,}0$)? Welcher Befund dominiert, und was bedeutet das für die Frage, ob Durchbrüche die Abflachungsrate zurücksetzen?}

\textbf{F3b~(TAS-Annäherungsanalyse)} zeigt, wie nah menschliche Weltrekorde dem theoretischen Minimum kommen und wie schnell die Lücke schließt. Der kombinierte Ansatz --- Asymptoten-Proxy für alle Spiele und direkter TAS-Vergleich für drei Spiele --- erlaubt eine Modellvalidierung und liefert spielspezifische Projektionen.

\placeholder{Kernergebnis von F3b eintragen: Mittlere Gap-Prozent; bei welchen Spielen ist die menschliche Community dem TAS am nächsten; wie gut trifft der Modell-Floor die echten TAS-Zeiten ($\Delta$)?}

\section*{Ausblick}

Weiterführende Analysen könnten die Stichprobe auf 100 oder mehr Spiele erweitern und zusätzliche Kategorien (100\%, Glitchless) einbeziehen. Eine Verknüpfung mit Community-Metriken (Anzahl aktiver Speedrunner, Twitch-Zuschauerzahlen) würde ermöglichen, die beobachteten Muster kausal zu erklären. Maschinelle Lernverfahren könnten eingesetzt werden, um die verbleibende Optimierungsspanne für aktive Rekorde vorherzusagen.
